\documentclass[10pt]{article}
\usepackage[margin=20mm]{geometry}
\usepackage{amsmath,amssymb,mathtools,fontspec,unicode-math,xurl,hyperref}
\IfFontExistsTF{DejaVu Serif}{\setmainfont{DejaVu Serif}}{\setmainfont{Latin Modern Roman}}
\setmathfont{latinmodern-math.otf}
\setlength{\parindent}{0pt}\setlength{\parskip}{6pt}\sloppy
\title{Technical Verification of
Sharp Bounds for Max-sliced Wasserstein Distances}
\author{AI-assisted technical review}
\date{}
\begin{document}
\maketitle
Abstract: This initial review concerns arXiv:2403.00666v6, not an unverified equivalent of the published article. The empirical-process argument, the symmetric quadratic lifting, and the lower-bound construction were examined. One limited sample-size defect is established: the literal logarithmic bound in Theorem 1.4 fails at \(n=1\). An exact counterexample, its operative hypotheses, the affected proof step, and a repair are given. This does not contradict the asserted large-sample rates. No major defect is established within the checked scope. Scientific adjudication remains pending.\par
\section*{Paper identity and scope}
The reviewed paper is March T. Boedihardjo, Sharp Bounds for Max-sliced Wasserstein Distances, arXiv:2403.00666v6, dated 26 November 2024, 27 PDF pages [source]. It is associated with DOI 10.1007/s10208-025-09690-1 (https://doi.org/10.1007/s10208-025-09690-1) in Foundations of Computational Mathematics. The publisher landing page lists online publication on 22 January 2025 and volume 26, pages 747–778 (2026) [journal]. That page supplied metadata and a subscription preview; the journal’s complete theorem text was not obtained. This report therefore attributes every mathematical finding to the acquired v6 PDF. Printed page numbers and one-based PDF page numbers coincide.\par
The complete v6 text was read, and all 27 rendered source pages were visually inspected in this scientific conversation. The review followed Theorems 1.1, 1.2 and 1.4, Corollary 1.3, Proposition 1.5 and the supporting arguments in Sections 2–3. It checked the applicability, rather than supplying new proofs, of the invoked majorizing-measure, Dudley, symmetrization and Khintchine inequalities. The external references were not all retrieved or independently audited. The independent computation below uses newly written rational-arithmetic code; no author code was executed. Missing optional code is not treated as a mathematical defect.\par
A bounded correction/version check consulted the official arXiv record and the publisher page, together with a title-and-correction search. The arXiv record still listed v6 as its latest version. No relevant author correction was identified in these consulted records. This is not evidence that no correction exists elsewhere, nor that the published text retains the issue.\par
\section*{Definitions and assumptions retained}
For probability measures on Euclidean space, the paper’s max-sliced distance is \[W_{p,1}(\mu,\nu)=
 \sup_{\|v\|_2=1}W_p\bigl((\langle v,\cdot\rangle)_\#\mu,
                         (\langle v,\cdot\rangle)_\#\nu\bigr).\] The more general \(W_{1,s}\) of Section 1.3 uses \(s\) coordinate functionals in the dual unit ball; their mutual orthogonality is not an assumption. In the Banach-space results, separability of the dual and the first-moment assumption are retained. In the Hilbert-space moment bounds, the stated \(2+\delta\) or \(2+2\delta\) moments are required where they occur, with \(0<\delta\leq1\). The centered lower bounds require \(\mathbb E X=0\); centering is not silently imposed on the upper bounds.\par
Equation reference: eq:original. Theorem 1.4 assumes a symmetric probability measure on \(\{x\in\mathbb R^d:\|x\|_2\leq r\}\), where \(r>0\), and independent samples \(X_1,\ldots,X_n\) with that law. Write \[\nu_n=\frac1{2n}\sum_{i=1}^n(\delta_{X_i}+\delta_{-X_i}),
 \qquad
 \Sigma=\mathbb E XX^{T},\qquad \kappa=\|\Sigma\|_{\rm op}.\] For \(\kappa>0\), its asserted upper bound is equivalently \[
 \mathbb E W_{2,1}(\mu,\nu_n)^2
 \leq C\left(\frac{r^2\log n}{n}
       +r\sqrt{\frac{\kappa\log n}{n}}\right).\] Neither copy of the theorem, on pages 6 and 22, explicitly states \(n\geq2\). The counterexample below uses \(\kappa>0\) and does not exploit division by zero or an inadmissible distribution.\par
\section*{A sample-size boundary defect}
Problem 1 (Theorem 1.4 at \(n=1\); minor severity).  The literal all-positive-sample-size reading of [eq:original] (\#eq:original) is false. The same boundary is absent from Lemma 3.1, on pages 20–21, which supplies the matrix estimate used in its proof.\par
\subsection*{Exact input and decisive calculation}
Equation reference: eq:mu. Set \(d=2\), \(r=1\), \(n=1\), and \[
 \mu=\tfrac14(\delta_{e_1}+\delta_{-e_1}
                  +\delta_{e_2}+\delta_{-e_2}).\] This is a probability measure, is symmetric, has support in the closed unit ball, and has covariance \(\Sigma=I_2/2\), so \(\kappa=1/2\). A single draw satisfies the independence requirement vacuously.\par
Claim 1.  For every possible value of \(X_1\) in [eq:mu] (\#eq:mu), \[W_{2,1}\bigl(\mu,\tfrac12(\delta_{X_1}+\delta_{-X_1})\bigr)^2
 =\tfrac12 .\]\par
Equation reference: eq:direction. Proof. Let \(v=(a,b)\) with \(a^2+b^2=1\), and put \(A=|a|\), \(B=|b|\). The absolute-value image of the projected law of \(\mu\) is \(\tfrac12\delta_A+\tfrac12\delta_B\). If the sampled axis is the first, the corresponding image of the symmetrized empirical law is \(\delta_A\); if it is the second, it is \(\delta_B\). For symmetric one-dimensional laws, quadratic transport equals quadratic transport of their absolute-value images. This follows from monotone quantiles, or by optimally coupling the nonnegative radii and assigning the same independent sign to both variables. The opposite inequality follows from the fact that absolute value is 1-Lipschitz. Thus in either case the projected squared distance is \[
 \tfrac12(A-B)^2=\tfrac12(1-2|ab|)\leq\tfrac12.\] Equality holds at \(v=e_1\) or \(v=e_2\). Taking the supremum proves the claim, including the zero-coordinate cases. ◻\par
Consequently the left side of [eq:original] (\#eq:original) is \(1/2\), whereas the right side is zero because \(\log1=0\), for every finite choice of the universal constant \(C\). This disproves the displayed endpoint assertion itself, not merely one method of proving it.\par
An independent exact-rational implementation of one-dimensional monotone transport checked 48 cases: all four sample points at 12 rational unit directions. Every computed value agrees with [eq:direction] (\#eq:direction); both coordinate directions attain \(1/2\). The saved script, output and execution receipt are `verification/0441/checks/boundary\_n1.py`, `boundary\_n1.stdout.json` and `execution.json`. The universal statement is proved above; the finite check is corroboration.\par
\subsection*{Proof dependency and alternative interpretation}
On page 20, Lemma 3.1 chooses the Schatten moment \(p=\lceil\log n\rceil\). At \(n=1\), this gives \(p=0\), outside the positive-moment argument, and the resulting Gaussian matrix estimate would have a zero right side. For example, \(\mathbb E\|gXX^T\|_{\rm op}=\mathbb E|g|>0\) when \(\|X\|_2=1\). Even the lemma’s first inequality fails literally at \(n=1\): uniform mass on \(\{\pm e_1,\pm e_2,\pm e_3\}\) gives \(\mathbb E\|XX^T\|_{\rm op}=1\) and \(2\|\mathbb E XX^T\|_{\rm op}=2/3\).\par
The dependency is direct: Lemma 3.2 and equation (3.2) lift the symmetric transport problem to the matrix law of \(XX^T\); Corollary 2.9 supplies the empirical-process bound; Lemma 3.1 then yields the logarithmic estimate on page 23. Corollary 2.9 also selects \(\delta=1/\lceil\log n\rceil\) on page 18. Its conclusion, however, need not fail at \(n=1\): \[\mathbb E W_{1,1}(\mu,\delta_X)\leq2\mathbb E\|X\|,
 \qquad
 \mathbb E\|gX\|=\sqrt{2/\pi}\,\mathbb E\|X\|,\] so its Gaussian term already controls that endpoint after an absolute adjustment of \(C\). This report does not label Theorem 1.2 or Corollary 2.9 false on that account.\par
An intended restriction \(n\geq2\) is plausible and is supported by the two ceiling-log choices in the proofs. Under that explicit restriction, the counterexample is outside the theorem’s scope. No global restriction or logarithm-truncation convention resolving the literal wording was identified in the reviewed v6 text. This interpretative boundary is why the finding is classified as minor: it changes a stated finite-sample endpoint, while the nontrivial asymptotic claim and its \(n\geq2\) argument survive.\par
\subsection*{Repair and impact}
Equation reference: eq:repair. One sufficient repair is to state \(n\geq2\) explicitly in the logarithmic matrix lemma and Theorem 1.4, and handle \(n=1\) separately where Corollary 2.9 is used. Alternatively, let \(L(n)=\log(en)\) and formulate the upper bound as \[
 \mathbb E W_{2,1}(\mu,\nu_n)^2
 \leq C\left(\frac{r^2L(n)}n+
                 r\sqrt{\frac{\kappa L(n)}n}\right),\qquad n\geq1.\] For \(n\geq2\), this follows from the checked argument for [eq:original] (\#eq:original), because \(L(n)\geq\log n\). For \(n=1\), both projected laws are supported on \([-r,r]\), so their squared quadratic distance is at most \(4r^2\), uniformly in direction and sample. Since \(L(1)=1\), [eq:repair] (\#eq:repair) follows after enlarging \(C\) to at least 4. This is a proved endpoint extension, not a claim that the author adopted it. Writing the bound in the expanded form also avoids undefined ratios when \(\kappa=0\); then \(\mu=\delta_0\) and the distance is zero.\par
\section*{Other checked arguments and counterevidence}
The following checks delimit the finding and record why nearby claims were not rejected.\par
\subsection*{Empirical-process upper bounds}
The product-metric estimate of Lemma 2.1, the weighted Lipschitz covering argument in Lemmas 2.3–2.4, and the Gaussian comparison in Lemma 2.5 were followed on pages 8–16. In particular the dyadic shell sum is finite for \(\delta>0\); the truncation parameter produces the distinct \(s=1\), \(s=2\), and \(s\geq3\) regimes in (2.12). Rescaling in Lemma 2.6 preserves both sides of (2.13); the all-zero sample is handled directly. Gaussian symmetrization in Theorem 2.7 applies to the centered Lipschitz class with the first-moment envelope. The Hilbert second-moment identity and Jensen’s inequality in Corollary 2.8 give the stated moment dependence. No contradictory example or invalid operative step was established in these arguments.\par
For Corollary 1.3, bounded support gives \(W_p^p\leq(2r)^{p-1}W_1\) after projection. Taking the supremum and then using concavity of the \(p\)-th root yields the claimed rate from Theorem 1.1. These steps retain \(p\geq1\). They do not require the additional symmetry assumption of Theorem 1.4.\par
\subsection*{Centered lower bounds and the covariance construction}
For Proposition 2.10, the linear test function supplies the norm of the empirical mean. Symmetrizing with an independent centered copy gives the factor \(1/2\), and conditional Jensen removes that copy. The Hilbert-space Khintchine step in Corollary 2.11 then gives the stated \(1/(2\sqrt{2n})\) coefficient. The zero-mean hypothesis is used essentially.\par
In Lemma 3.2, symmetry permits passage to absolute values, and \((u-v)^2\leq|u^2-v^2|\) for \(u,v\geq0\) permits the square-map comparison. The trace functional represented by \(ww^T\) has dual norm one when \(\|w\|_2=1\). The argument would not justify dropping symmetry; the reviewed theorem does not drop it.\par
For Proposition 1.5, normalize \(\operatorname{Tr}\Sigma=1\) in the nonzero case and write its eigenvalues as \(\lambda_j\), with \(\lambda_1\leq1/2\). The symmetric coordinate-axis law has the required covariance. The occupancy-vector computation on page 24 gives \[\mathbb E W_{2,1}(\mu,\nu_n)^2
 \geq \frac{1-\sum_j\lambda_j^2}{4n}\geq\frac1{8n}.\] For the complementary regime \(n\lambda_1\geq1\), projection onto the first axis gives \(\mathbb E|Y/n-\lambda_1|\), where \(Y\) is binomial. As an independent sufficiency check, put \(Z=Y-np\), \(p=\lambda_1\), and \(\sigma^2=np(1-p)\geq1/2\). Direct moment expansion gives \(\mathbb E Z^4=3\sigma^4+\sigma^2(1-6p(1-p))
\leq5\sigma^4\). Hölder’s inequality implies \[\mathbb E|Z|\geq
 \frac{(\mathbb E Z^2)^{3/2}}{(\mathbb E Z^4)^{1/2}}
 \geq\frac{\sigma}{\sqrt5},
 \qquad
 \mathbb E|Y/n-p|\geq\sqrt{\frac{p}{10n}}.\] Together with \(1/(8n)\), this is already sufficient for the paper’s \(1/16\) lower bound, by considering whether \(\sqrt{p/n}\) is greater than \(1/n\). It does not verify the sharper external binomial inequality cited on page 25; it independently confirms that the main constant claimed here does not depend on that sharper value.\par
\section*{Conclusion}
The established defect is confined to the literal \(n=1\) scope of the logarithmic covariance estimate and Theorem 1.4. The counterexample is exact and satisfies every hypothesis listed in those statements. An explicit \(n\geq2\) restriction, or the proved extension [eq:repair] (\#eq:repair), resolves it without altering the reported large-sample orders. No major or critical issue is established by this review. This is a bounded initial technical report, not an independent acceptance decision.\par
Assessment | Repair verification | Author adoption | Remaining scope\par
Minor | \(n\geq2\) restriction; proved \(\log(en)\) extension | Unknown | Journal full text not reviewed; independent acceptance pending\par
\section*{Evidence and review limits}
The source PDF has SHA256\par
`8b5e27e131dfda0c72e5983d588b5be6374c9`
`b5e4c6032d01f25abf8ea3a6287`.\par
The manifest, exact source, independent check, source-page inspection receipt, build records, final PDF inspection facts and support inventory are retained under the matching `0441` trees in the assigned stage. Hashes and execution facts in those records bind the actual artifacts; this prose does not substitute for a build or a viewing event.\par
All substantive reading and reasoning reported here were performed in the assigned scientific conversation. There was no independent human review and no author-code execution. The management UI observation and the conversation’s runtime self-description are recorded separately in provenance; no independently verified backend identity is asserted. Scientific decision: pending. Human review records: empty.\par
9 M. T. Boedihardjo, Sharp Bounds for Max-sliced Wasserstein Distances, arXiv:2403.00666v6, 26 November 2024, 27 pp. https://arxiv.org/abs/2403.00666v6. Exact acquired PDF reviewed; official version history also consulted. Springer Nature, article landing page for M. T. Boedihardjo, Sharp Bounds for Max-sliced Wasserstein Distances, Found. Comput. Math. 26 (2026), 747–778, online 22 January 2025. https://doi.org/10.1007/s10208-025-09690-1. Metadata/subscription preview only; complete journal text not reviewed.\par
\end{document}
